%=====================================================================
% 第 7 节  对 sigma 模型的诠释？
% 原文：Section 7. Sigma model interpretation?（arXiv:1702.06499v2）
%=====================================================================
\section{对 sigma 模型的诠释？}

\begin{keytakeaways}
\begin{itemize}
  \item \textbf{物理图像：}把 AGT 与几何 Langlands 的图像放回到以 $\CM_H(G)$ 为目标流形的
        二维拓扑 sigma 模型中。Kapustin--Witten 用的是以 $\omega_K$ 为辛结构的 A-模型、
        其 cc-brane 的 Chan--Paton 曲率为 $\omega_J$；Nekrasov--Witten 用的是以
        $\frac{\ep_2}{\ep_1}\omega_I$ 为辛结构的 A-模型。二者之间没有明显参数可以连续移动，
        这正是本节要化解的障碍。
  \item \textbf{数学工具：}Hitchin 模空间的超 Kähler 旋转（Hitchin 的 $\BC^*$ 作用把
        $(I_\zeta,\Omega_\zeta)$ 族中除 $\zeta=0,\infty$ 外的复结构与辛结构互相联系起来）、
        SYZ 镜对称与 Hitchin 纤维化、非阿贝尔 Hodge（NAH）对应及其共形极限、以及
        $\ep_2\ra0$ 下的 WKB 近似。
  \item \textbf{关键结论：}若 sigma 模型 TQFT 在 $\BC^*$ 作用下具有适当形式的不变性，则取
        $\zeta=\ep_2\xi$ 即可把 NW 与 KW 的两个 A-模型插值起来；在 NS 极限下 brane
        $\mathsf{L}_{\si,a}^{\sst(1)\hbar}$ 退化为支撑在 Hitchin 第一纤维化纤维
        $\mathsf{F}_{(u_{\mathsf{a}},0)}$ 上、Chan--Paton 丛平凡的 (B,A,A)-brane，于是
        $\ep_2\ra0$ 时共形块与 Kapustin--Witten 的 Hecke 本征 brane 之间的对应可写成显式同构。
  \item \textbf{技术要点：}该同构的自然形式是把 Beilinson--Drinfeld 进路中的 oper
        $\rho_u$ 经 NAH 对应的共形极限（$R\ra0$、$\zeta\ra0$、$\zeta/R$ 有限）送到 Hitchin
        模空间的 $(u,0)$ 点；对一般局部系 $\chi$ 则需假设 $\CM_{dR}({}^LG)$ 与
        $\CM_H({}^LG)$ 之间存在由该共形极限相联的分层。
\end{itemize}
\end{keytakeaways}

\bigskip
\noindent\rule{\textwidth}{0.6pt}
\bigskip

上一节的讨论揭示了（面算符存在时）AGT 对应的各种推广与几何 Langlands 对应的
Beilinson--Drinfeld 进路之间的密切联系。我们已经把几何 Langlands 对应识别为：
用“量子化的特征标簇”描述共形块空间这一描述在 $\ep_2\ra 0$ 时的极限。
在这最后一节中，我们要回到这样一个问题：由此浮现的图像与 Kapustin--Witten 进路下的
几何 Langlands 对应之间是什么关系。在 \ref{comp-I} 小节中我们提出如何在 $2d$ sigma 模型中
描述极限 $\ep_2\ra 0$ 的问题；\ref{Cstaraction} 与 \ref{SecNSlim} 两小节给出这一问题的一个
推测性回答。本节其余部分将基于类似的论证，对 Beilinson--Drinfeld 与 Kapustin--Witten
两种进路作更明确的比较。

\subsection{比较 I：AGT 对应与几何 Langlands 对应之间的关系}\label{comp-I}

我们曾论证：此语境中出现的 $\CD$-模可以用相应微分方程的解空间来刻画，而这些解在物理上
实现为配分函数。在带面算符的 AGT 对应语境中，我们曾论证空间
$\CH_{x}^{\sst(2)}:=
{\rm Hom}_{\CM_H(G)}^{}(\mathsf{B}_{\rm cc}^{},\mathsf{L}_{x}^{\sst(2)})$
可以等同于仿射代数共形块在丛 $\CE_x$ 上的纤维。在 CFT 中，我们可以把解 KZB 方程的
手征配分函数 $\CZ(a;x,\tau;k)$ 看作共形块空间 $\mathcal{CB}(C,\hfg_k,\CE_x)$ 的子空间
$\mathcal{CB}^\si_a(C,\hfg_k,\CE_x)$ 中元素 $\vf_a$ 的代表元，而这些子空间由这样的共形块
组成：它们在 Verlinde 线算符的某个极大对易子集下具有由 $a$ 参数化的固定本征值。

开 sigma 模型 TQFT 的模不变性提示我们可以等同
$\mathcal{CB}^\si_a(C,\hfg_k)\simeq 
{\rm Hom}_{\CM_H(G)}^{}(\mathsf{B}_{\rm cc}^{},\mathsf{L}_{\si,a}^{\sst(1)})$。
这一等同可以由下述观察支持：边界条件 $\mathsf{L}_{\si,a}^{\sst(1)}$ 固定了条带上
sigma 模型中某些零模的取值，而这应当导致相对于 $\mathcal{A}_{B}^{\hbar}$ 的某个子代数的
本征值性质，该子代数由相应的量子化迹函数生成，而这些迹函数则与决定裤分解 $\si$ 的
那些曲线相联系。

在上一节中我们讨论了子空间 $\mathcal{CB}_a^\si(C,\hfg_k,\CE_x)$ 在极限 $\ep_2\ra 0$ 下如何
与几何 Langlands 对应中出现的 Hecke 本征 $\CD$-模相联系，其中代表“本征值”的局部系由
以 $a$ 为参数的 opers 族表示。然而不清楚的是：我们关于 Hecke 本征值性质在极限
$\ep_2\ra 0$ 下出现的观察，如何从 sigma 模型的角度来理解。似乎存在一个直接的障碍：
\cite{KW} 在几何 Langlands 对应语境中考虑的 sigma 模型，是以辛结构 $\omega_K$ 的 A-模型，
而 cc-brane 的 Chan--Paton 曲率为 $\omega_J$。在 AGT 对应的 Nekrasov--Witten 进路中则做了
不同的选择，那里用的是以辛结构 $\frac{\ep_2}{\ep_1}\omega_I$ 的 A-模型。没有任何明显的参数
能使我们在这两种情形之间连续移动。

鉴于在 CFT 中可以在临界水平极限下得到几何 Langlands 对应，很自然地要问：这能否在以
$\CM_H(G)$ 为目标流形的 $2d$ sigma 模型内理解。从 \ref{Cstaraction} 小节起，我们将推测
一条可能的途径。

\subsection{超 Kähler 旋转}\label{Cstaraction}

一个诱人的做法是：用来自 Hitchin 模空间上圆作用 $\vf\ra e^{i\theta}\vf$ 的超
Kähler$^{(g)}$ 旋转，来修改 Nekrasov 与 Witten 的二维 TQFT 构架。\cite[Section 9]{Hi87}
中证明：这一圆作用可以复化为一个 $\BC^*$-作用，它把族 $(I_\zeta,\Omega_\zeta)$ 中除
$\zeta=0,\infty$ 之外的所有复结构与辛结构互相联系起来。超 Kähler 旋转使我们能够在
$\zeta\neq 0,\infty$ 时识别 $\CM_H(G)$ 上的关键几何结构，其含义是这些旋转可以由
$\CM_H(G)$ 的微分同胚实现 \cite[Proposition (9.1)]{Hi87}。这相当于对 sigma 模型场作
重新参数化，而后者也许并不影响物理内容。

sigma 模型路径积分在场重定义下的不变性可能是个微妙的问题。圆作用
$\vf\ra e^{i\theta}\vf$ 应当容易理解；但复化为 $\BC^*$-作用之后，它是否在适当意义上
容易理解为 sigma 模型 TQFT 的对称性，我们并不清楚。我们想指出的是：若以 $\CM_H(G)$
为目标流形的拓扑 sigma 模型在 $\BC^*$-作用下具有某种适当形式的不变性，我们就能简单地
在 \cite{KW} 与 \cite{NW} 中分别考虑的那两个 sigma 模型之间插值。这样就能对上一节的观察
给出一个有吸引力的解释。

我们关心的是带边曲面上的开 A-模型，其边界条件或属典范余迷向（canonical coisotropic）型，
或属拉格朗日型。我们可以用全纯辛形式族 $\Omega_{\zeta}$ 来定义一族辛结构为
$\mathrm{Im}(\Omega_\zeta)$ 的 A-模型，其中余迷向 brane 的 Chan--Paton 曲率为
$\mathrm{Re}(\Omega_\zeta)$。$\CM_H(G)$ 中对 $\Omega_\zeta$ 为拉格朗日的子流形，在这族
A-模型中定义 brane。若 sigma 模型 TQFT 具有适当形式的 $\BC^*$-不变性，则可以预期：
在这族 A-模型中定义的配分函数 $\CZ_\zeta$ 实际上与 $\zeta$ 无关。

然而应当注意：我们用来定义边界条件的所有拉格朗日子流形都必须协调地变动，才能保持
$\CZ_{\zeta}$ 不变。我们所关心的子流形是特征标簇中复 Fenchel--Nielsen$^{(g)}$ 扭转流的
轨道，经由 NAH 对应与和乐映射视作 $\CM_H(G)$ 的子流形。固定特征标簇中的一条轨道，
随着超 Kähler 参数 $\zeta$ 的变动，就得到 $\CM_H(G)$ 的一族单参数子流形。

若已有 $\BC^*$-作用下的不变性，我们就可以把耦合参数 $\ep_1/\ep_2$ 的变动与适当的超 Kähler
旋转结合起来，从而得到一族以 $\CM_H(G)$ 为目标流形、并在 \cite{NW} 与 \cite{KW} 所考虑的
两类模型之间插值的拓扑 A-模型。这可以通过取 $\zeta=\ep_2\xi$ 来实现。表达式
$\hat{\omega}_\zeta=\frac{\ep_2}{\ep_1}\mathrm{Im}(\Omega_{\zeta})$ 在极限
$\ep_2\ra 0$ 下约化为 $\hat{\omega}_\zeta=
\frac{1}{2\ep_1|\xi|^2}
(\mathrm{Re}(\zeta)\omega_K-\mathrm{Im}(\zeta)\omega_J)$，
对实 $\xi$ 恰好重现用来定义 \cite{KW} 所研究的 A-模型的辛形式。

\subsection{brane $\mathsf{L}_{\si,a}^{\sst (1)}$ 的 Nekrasov--Shatashvili 极限} \label{SecNSlim}

为看清由此得到的方案是否真的能解决上述困惑，先注意：\ref{branes1} 小节所考虑的 (A,B,A)
brane $\mathsf{L}_{\si,a}^{\sst (1)}$ 可以很容易地推广为 $(B,A,A)_\zeta$-brane
$\mathsf{L}_{\si,a}^{\sst (1)\zeta}$，只要在它们的定义中使用依赖于 $\zeta$ 的 NAH 对应。
如上面建议的那样以依赖于 $\hbar$ 的方式选取 $\zeta$，就会把它们变成超 Kähler 参数为
$\hbar=\ep_1/\ep_2$ 的那族 A-模型中的 brane $\mathsf{L}_{\si,a}^{\sst (1)\hbar}$。

我们会看到，在 brane 族 $\mathsf{L}_{\si,a}^{\sst (1)\hbar}$ 的 Nekrasov--Shatashvili 极限中
发生了有趣的事情：若在此极限中把参数 $a$ 按
${a}_r=\frac{\ep_1}{\ep_2}\,\check{a}_r$ 标度（$\check{a}_r$ 有限），则 brane
$\mathsf{L}_{\si,a}^{\sst(1)\hbar}$ 有良定义的极限 $\hbar\ra \infty$。为理解极限
$\ep_2\ra 0$，主要需要研究 $\pa_u+\mathcal{A}_u$ 的和乐的 WKB 近似，其中
$\mathcal{A}_u=A_u+\frac{1}{\ep_2}\vf$。通过规范变换总可以局部达到形式
$\mathcal{A}_u=\big(\begin{smallmatrix} 0 & u_{\ep_2} \\ 1 & 0\end{smallmatrix}\big)$。
出现在 $\mathcal{A}_u$ 右上矩阵元中的函数 $u_{\ep_2}(t)$ 一般只是 $u$ 的亚纯函数，
但其留量为 $\ep_2$ 阶。我们注意到 $A_u$ 只出现在展开的次领头阶。由这些事实立即得到
$\vartheta_{\ep_2}(t)=\frac{1}{2}{\rm tr}(\vf^2)+\CO(\ep_2)$。因此在 $\ep_2$ 的领头阶，
可以把 $(\pa_u+\mathcal{A}_u)s_{\ep_2}=0$ 的解 $s$ 写成
\begin{equation}
s_{\ep_2}(t)\,\sim\, e^{\pm\frac{1}{\ep_2}\int_{\CC_t} \la}\,,
\end{equation}
其中 $\CC_t$ 是谱曲线 $\Sigma$ 上终止于 $t\in \C$ 的某个提升 $\hat{t}\in\Sigma$ 的路径，
而 $\la$ 是 $\Sigma$ 上的典范微分。结论是：标度后的 Fenchel--Nielsen 长度坐标
$\check{a}_r$ 在极限 $\ep_2\ra 0$ 下表现为
$
\check{a}_r=\frac{1}{\ep_1}\mathsf{a}_r+\CO(\ep_2),
$
其中 $\mathsf{a}_r$ 是微分 $\la$ 沿一组由 $\si$ 确定的典范同调基的某个拉格朗日子空间所取的
周期。

在这个意义上，人们可以把“量子周期”$\check{a}_r$ 看作 Hitchin 可积系统的角变量
$\mathsf{a}_r$ 的形变。有趣的是：对 $\si$ 的依赖在极限 ${\ep}\ra 0$ 下消失——即使坐标
$\mathsf{a}_r$ 的定义带有对 $\si$ 的残余依赖，由 $\mathsf{a}_r$ 的取值所确定的 Hitchin
纤维化的纤维却并非如此。

这样我们便被引向如下结论：brane $\mathsf{L}_{\si,a}^{\sst(1)\hbar}$ 的标度极限由
(B,A,A)-brane 表示，后者支撑在 Hitchin（第一）纤维化的纤维
$\mathsf{F}_{\sst (u_{\mathsf{a}},0)}$ 上，其上的 Chan--Paton 丛是平凡丛，而
$u_\mathsf{a}$ 是 Hitchin 基 $\CB$ 上由坐标 $\mathsf{a}_r$ 确定的点。

最后要注意：也可以用类似方式讨论 brane $\mathsf{L}_{x}^{\sst(2)\hbar}$。容易看出，
当 ${\hbar}\ra 0$ 时得到支撑在 Hitchin 第二纤维化的纤维上的 brane
$\mathsf{L}_{x}^{\sst(2)}$。

\subsection{比较 II：Kapustin--Witten 进路与 CFT 的关系？}\label{compII}

任何在 Beilinson--Drinfeld 与 Kapustin--Witten 两种进路之间的比较，都必须回应下面这一点。
在 \cite{KW} 中，Hecke 本征 brane 被描述为 $\CM_H(G)$ 上的摩天楼层（skyscraper sheaves）。
然而，为了在 Hitchin 纤维化上使用 SYZ 镜对称，Kapustin 与 Witten 采用的是适配复结构 $I$ 的
$\CM_H(G)$ 与 $\CM_H({}^L G)$ 之表示，而 Beilinson--Drinfeld 进路则在对应的一侧考虑
${}^LG$-局部系 $(\CE,\nabla')$。为把二者联系起来，需要用到 NAH 对应。若 $\mu$ 表示
$\CM_H(G)$ 的一个点（该空间表为环面纤维化），令 $\check{\mathsf{F}}_{\mu}$ 为 B-模型中
相应的摩天楼层，$\mathsf{F}_{\mu}$ 为 SYZ-对偶的 A-模型中的 Hitchin 纤维。为表述两种进路
之间的一个猜想性关系，我们需要用 NAH 对应（的逆）找出与局部系 $\chi$ 相联系的点
$\mu(\chi)\in\CM_H(G)$。

于是，对 Beilinson--Drinfeld 与 Kapustin--Witten 两种进路之间可能关系的一个自然猜测，
就是下述 $\CD$-模同构成立：
\begin{equation} \label{BD-KW}
\begin{aligned}
{\rm Hom}_{\CM_H(G)}^{}(\mathsf{B}_{\rm cc},\mathsf{F}_{\mu(\chi)}) 
& \;\simeq\; \mathcal{CB}(C,\hfg_{-h^{\vee}},\CE)\;\;\text{for}\;\;\CE\in\mathrm{Bun}_G^{\rm vs}\\
& \;\simeq\; \big[\,
\text{Fiber of}\;\, \De_\chi\;\,\text{over}\;\,\CE\in\mathrm{Bun}_G^{\rm vs}\,\big]
\end{aligned}
\end{equation}
其中 $\mathrm{Bun}_G^{\rm vs}$ 是“非常稳定”（very stable）丛 $\CE$ 的空间，即不容许幂零
Higgs 场的丛；$\De_\chi$ 是由 $\mathcal{CB}(C,\hfg_{-h^{\vee}},\CE)$ 表示的 $\CD$-模。
由于存在一个典范平坦联络把与不同 $\CE\in\mathrm{Bun}_G^{\rm vs}$ 相联系的纤维等同起来，
右端不依赖于 $\CE$ 的选取。$\mathcal{CB}(C,\hfg_{-h^{\vee}},\CE)$ 的维数在离开
$\mathrm{Bun}_G^{\rm vs}$ 时可能发生跳跃，讨论见 \cite[Section 9.5]{Fr07}。

在此我们还可以指出 Beilinson--Drinfeld 与 Kapustin--Witten 进路比较中出现的另一个困惑。
如上文所指出的、并由 \cite{T10} 所研究的例子所说明的：在 Beilinson--Drinfeld 进路中出现的
$\CD$-模，在离开 $\CM_{dR}({}^LG)$ 内的 opers 子簇时表现出某种不连续的行为，其表现是：
在这些 $\CD$-模以共形块形式给出的表示中，会出现若干个额外的退化表示。而在 Kapustin 与
Witten 的进路中看不到这种不连续行为。

\subsection{共形极限}

我们还愿意指出一条也许能为 \ref{compII} 小节提出的问题提供答案的途径。它基于
\cite{Ga14,DFKMMN} 中的观察：在共形极限下 NAH 对应可能大大简化；该极限是把通过
$\vf\ra R\vf$ 引入 NAH 对应的参数 $R$ 与超 Kähler 参数 $\zeta$ 一起趋于零，同时保持
$\zeta/R$ 有限。作为准备，这里先讨论该极限可能具有的相关性。所涉及的问题与
\ref{Cstaraction} 小节中讨论的超 Kähler 旋转的相关性问题类似，但并不相同，因此值得
单独讨论。

作替换 $\vf\ra R\vf$ 定义了 NAH 对应的一族单参数形变，看起来会修改 $\CM_H(G)$ 上的
超 Kähler 度量。然而，参数 $R$ 对 Hitchin 模空间的几何而言是非本质的，其含义是：
不同 $R$ 值对应的超 Kähler 度量之间由微分同胚相联系。人们可以用由修改后的 NAH 对应
所得的 $\CM_H(G)$ 上的超 Kähler 度量，来定义一族单参数 sigma 模型作用量 $S_R$。
但是，既然这些作用量是由以微分同胚相联的超 Kähler 度量构造的，人们或会倾向于把仅在
参数 $R$ 的选取上不同的两个作用量 $S_{R_1}$ 与 $S_{R_2}$，认同为同一个 sigma 模型量子场论
之物理等价的拉格朗日表示。在诸如 $u=\frac{1}{2}\mathrm{tr}(\vf^2)$ 这样的合适坐标中，
可以发现相应的微分同胚实现为简单的标度变换 $u\ra ur^{2}$。相应的场重定义确实只应以
非本质的整体因子来重新标度配分函数。

若用作用量 $S_R$ 定义的配分函数 $Z_R(u)$ 依赖于一个边界参数 $u$，而该参数对应于
$\CM_H(G)$ 的一个坐标、在 $\vf\ra R\vf$ 下按 $u\ra ur^{2}$ 标度，则我们发现：sigma 模型
度量在微分同胚意义下的 $R$-无关性表现为 $Z_{r R}(u)\propto Z_R(ur^{2})$，同样可能只差
非本质的整体因子。

边界 B-模型中由 $\CM_{dR}({}^LG)$ 上的点 $\chi$ 定义的边界条件所对应的配分函数
$Z'(\chi)$，可以用依赖于点 $\mu_{R,\zeta}(\chi) \in \CM_{H}({}^LG)$ 的配分函数
$Z_R''(\mu_{R,\zeta}(\chi))$ 来表示，因为由 NAH 对应之修改所产生的 $R$-依赖，被 sigma 模型
作用量的相应修改所补偿。若存在使 $\mu_{R,\zeta}(\chi)$ 大大简化的极限，这一观察可能有用。
我们将看到，这样的极限正是所谓的共形极限 $R\ra 0$、$\zeta\ra 0$ 且保持 $\zeta/R$ 有限。
如下文所述，它将给出所提议的关系式\eqref{BD-KW} 的一个有趣的重述。

\subsection{更明确地表述 Kapustin--Witten 进路与 CFT 的关系}

我们将主要限于这样的情形：局部系 $\chi$ 为 opers
$\rho_u=\big(\CE_o,\big(\pa_t+\big(\begin{smallmatrix} 0& u\\ 1 & 0\end{smallmatrix}\big)\big)dt\big)$，
它们在 Beilinson 与 Drinfeld 的进路中起基本作用。有趣的是：在此情形下 NAH 对应的共形极限
变得特别简单 \cite{DFKMMN}，它把 opers $\rho_u$ 与一类形式非常特殊的 Higgs 对联系起来。
对 ${}^L G=\mathrm{PSL}(2,\BC)$，相应的 Higgs 对 $(\mathcal{E},\vf)$ 形如
$(K^{1/2}\oplus K^{-1/2},(\begin{smallmatrix} 0& u \\ 1 &0\end{smallmatrix})dz)$，其中 $u$
是在 $C$ 的 Fuchs 单值化中表示 $\rho_u$ 的那个二次微分。若把 $\CM_H({}^L G)$ 描述为由
$(\Sigma_u,\CL)$ 这样的对表示的环面纤维化（$\Sigma_u$ 是与 Hitchin 基上一点 $u\in\CB$ 相联系的谱曲线，
$\CL$ 是 $\Sigma_u$ 上的线丛），则得到线丛 $\CL_0=\pi^*(K^{1/2})$，其中 $\pi: \Sigma_u\ra C$
是覆盖投影。该线丛代表参数化 $\CL$ 各种选择的 Jacobian/Prym 的一个典范“原点” \cite{W15}。

令 $\check{\mathsf{F}}_{\sst(u,0)}$ 为 $\CM_{H}({}^LG)$ 上支撑在点 $(\Sigma_u,\CL_0)$ 的
摩天楼层，并令 $\mathsf{F}_{\sst (u,0)}$ 为与 $\check{\mathsf{F}}_{\sst(u,0)}$ SYZ 对偶的
Hitchin 纤维化的纤维。于是，Beilinson--Drinfeld 与 Kapustin--Witten 两种进路之间的可能
关系式\eqref{BD-KW} 可以更明确地写成
\begin{equation} \label{BD-KW-conf}
{\rm Hom}_{\CM_H(G)}^{}(\mathsf{B}_{\rm cc},\mathsf{F}_{\sst (u,0)})
\;\simeq\;
\big[\,
\text{Fiber of}\;\, \De_{\rho_u}\;\,\text{over}\;\,\CE\in\mathrm{Bun}_G^{\rm vs}\,\big]\,.
\end{equation}
我们指出：$\CM_{H}({}^LG)$ 中一般点 $(u,0)$ 在 NAH 对应下的像，当且仅当从
$\CM_{H}({}^LG)$ 到 $\CM_{dR}({}^LG)$ 的映射是用非阿贝尔 Hodge 对应的共形极限定义时，
才由某个 oper 表示。\footnote{“if”这一方向已在 \cite{DFKMMN} 中证明。
$(u,0)\in\CM_{H}({}^LG)$ 中的点在 NAH 下的像由具有实和乐的联络组成，这些联络只离散地与
opers 簇相交。我们感谢 A. Neitzke 的这一说明。} 这一点，以及该极限在 \ref{Hecke-4d}
小节中的相关性，表明这一极限很适合用来表述 Beilinson--Drinfeld 与 Kapustin--Witten
两种进路之间的关系。

关于把 \eqref{BD-KW-conf} 推广到更一般的局部系 $\chi$，我们猜想：存在
$\CM_{dR}({}^LG)$ 与 $\CM_{H}({}^LG)$ 的自然分层，其各层之间由 NAH 对应的共形极限
相联系。这将使我们能把关系式\eqref{BD-KW-conf} 推广到一般的不可约局部系，从而把几何
Langlands 对应中 $\CD$-模的不连续行为与从一个层过渡到另一个层联系起来。我们计划在后续
文章中回到这一点。

\subsection{展望}

我们将在别处讨论形如
\begin{equation}\label{spectral}
\CH_{x}^{\sst(2)}\simeq \int^{\oplus} d\mu_\si(a)\;\CH_{\si,a}^{\sst(1)}\,,
\end{equation}
的关系之存在的现有证据，以及关于与不同裤分解 $\si$ 相联系的各空间
$\CH_{\si,a}^{\sst(1)}$ 之间存在线性关系的证据。这就把较弱版本的 $\si$-无关性恢复到了
与非零 $\epsilon_2$ 相联系的那套图景之中。

几何 Langlands 对应有时被表述为经典 Langlands 纲领中出现的自守形式空间之谱分解的类比。
我们把本短文的内容视为一些提示：梦想一个由超越（transcendental）与分析（analytic）两方面
扩展的几何 Langlands 纲领的变体，或许并不算放肆。超越方面涉及配分函数，它们表示由
$\CD$-模所定义的微分方程组的解；分析方面则涉及如 \eqref{spectral} 所提议的谱分解。
把配分函数认同为自守形式的类比物，会进一步加强对经典 Langlands 纲领的类比。配分函数
构成了 $\CM_H(G)$ 的代数结构之间的桥梁，而这些代数结构分别与 $\CM_H(G)$ 作为局部系
模空间 $\CM_{dR}(G)$ 的表示、以及作为特征标簇 $\CM_B(G)$ 的表示相联系。由此可以期望得到
一个更大的图像，把几何 Langlands 纲领中拓扑的方面与依赖复结构的方面统一起来。

我们计划在后续文章中讨论这些内容、“量子几何 Langlands 对偶”这一解释，以及与 \cite{T10}
中称为模对偶（modular duality）的另一种 Langlands 对偶模式的化身之间的关系。

最后注意，从规范场论方面对几何 Langlands 纲领的近期进展已在 \cite{Ga16a,Ga16b} 中描述。
分析它们与我们工作的关系应当是有意义的。

{\bf 致谢：}
A.B 感谢 D. Ben-Zvi、M. Mulase、A. Neitzke 与 R. Wentworth 的讨论，并感谢 String-Math 2016
的组织者筹办了一场富有启发的会议。A.B 还感谢 L'Institut Henri Poincar\'e（巴黎）与
ICTS-TIFR（班加罗尔）在工作进行期间的接待。

J.T. 感谢 String-Math 2016 的组织者筹办了一场鼓舞人心的会议并提供报告本文工作的机会，
感谢 M. Mulase、A. Neitzke 的讨论。特别感谢 E. Frenkel 的多方讨论、告知其未发表工作
\cite{Fr10} 的内容，以及对草稿的批评性意见。

本工作由德国科学基金会（DFG）通过协作研究中心 SFB 676 “Particles, Strings and the Early
Universe” 项目 A10 资助。

\bigskip
\noindent\rule{\textwidth}{0.6pt}
% [Tutor Note: 超 Kähler 旋转这一步最易被误读。Hitchin 的 $\BC^*$-作用对实 $\theta$ 是 $U(1)$ 等距，但复化之后**不是等距**：\cite[Proposition (9.1)]{Hi87} 只断言 $\zeta\neq0,\infty$ 时超 Kähler 旋转由 $\CM_H(G)$ 的微分同胚实现，而该微分同胚会改变度量（它是“旋转复结构”而非“保持度量”）。因此“这相当于对 sigma 模型场作重新参数化，也许并不影响物理内容”是一个需要证明的断言，而非平凡事实：sigma 模型的作用量依赖目标度量，一般的目标空间微分同胚（非等距）会改变作用量；此外场重定义还会移动 $B$ 场与路径积分测度，除非同时作相应补偿且无反常，物理才真正不变。原文自己也承认“may be a subtle issue”。若要在 KW 与 NW 两个 A-模型之间插值，需要的正是这一不变性；而 A-模型只依赖辛形式，故真正要验证的是 $\mathrm{Im}(\Omega_\zeta)$（连同 $\mathrm{Re}(\Omega_\zeta)$ 作为 Chan--Paton 曲率）在 $\BC^*$-作用下的行为。]
% [Tutor Note: 第 7.2 小节末的 $\hat\omega_\zeta$ 极限式（原文未编号）请用附录中的公式复核。附录给出 $\Omega_\zeta=\frac{1}{2\zeta}(\omega_J+i\omega_K)+i\omega_I+\frac{\zeta}{2}(\omega_J-i\omega_K)$。代入 $\zeta=\ep_2\xi$ 得 $\frac{1}{\zeta}=\bar\xi/(\ep_2|\xi|^2)$，故 $\Omega_\zeta\simeq\frac{1}{2\ep_2|\xi|^2}\big[(\mathrm{Re}\xi\,\omega_J+\mathrm{Im}\xi\,\omega_K)+i(\mathrm{Re}\xi\,\omega_K-\mathrm{Im}\xi\,\omega_J)\big]+i\omega_I$，于是 $\mathrm{Im}\Omega_\zeta=\frac{1}{2\ep_2|\xi|^2}(\mathrm{Re}\zeta\,\omega_K-\mathrm{Im}\zeta\,\omega_J)+\omega_I$。乘上 $\ep_2/\ep_1$ 后，$\omega_I$ 项给出 $\frac{\ep_2}{\ep_1}\omega_I$，它在 $\ep_2\ra0$ 时被丢掉——原文未曾明言这一丢弃，而这一项正是“NW 的 $\omega_I$ 模型”与“KW 的 $\omega_K$ 模型”之间差别的残留。另注意对实 $\xi$ 只得到正比于 $\omega_K$ 的辛形式，整体因子为 $1/(2\ep_1\xi)$，故与 \cite{KW} 的 $\omega_K$ 只相差一个常数标度；常数标度在 A-模型中只改变整体因子，不影响 brane 的拉格朗日性。]
% [Tutor Note: WKB 与标度这一步有三处易错。(1) 指数里是 $\la$ 而不是 $u$ 或 $\sqrt{u}$：$\la$ 是谱曲线 $\Sigma\subset T^*C$ 上的典范微分，$\pm$ 对应两叶，$\CC_t$ 是到提升点 $\hat t$ 的路径；$1/\ep_2$ 来自 $\mathcal{A}_u$ 中 $\vf$ 前的 $1/\ep_2$，故 WKB 的大参数是 $\hbar=\ep_1/\ep_2$（在 $\ep_2\ra0$、$\ep_1$ 固定时趋于无穷），不要误写成 $1/\ep_1$。(2) 这是奇异极限，$e^{\pm\frac{1}{\ep_2}\int\la}$ 只在各 Stokes 区域内渐近有效，绕行转折点（$\la=0$ 处）时按 Stokes 矩阵跳跃，真正的解须经 Borel 求和才确定；原文直接写“解 $s$ 的形式”跳过了 Stokes 现象，而 \cite{GMN2} 表明这一跳跃正对应谱曲线的墙穿越（WKB 曲线的 BPS 态计数），它会影响下文 $\check a_r$ 的 $\CO(\ep_2)$ 修正。(3) 标度关系可疑：把 $a_r=\frac{\ep_1}{\ep_2}\check a_r$ 与 $\check a_r=\frac{1}{\ep_1}\mathsf a_r+\CO(\ep_2)$ 联立得 $a_r=\frac{1}{\ep_2}\mathsf a_r+\CO(\ep_1)$，令 $\ep_2\ra0$ 会发散；而若 $\mathsf a_r$ 确如正文所说是 $\la$ 的周期、且经典长度坐标即该周期，则更自然的写法应为 $\check a_r=\frac{\ep_2}{\ep_1}\mathsf a_r+\CO(\ep_2)$。原文此处 $\ep_1,\ep_2$ 的归属值得核对（译文保留原样）。此外下一句把 $\mathsf a_r$ 称为 Hitchin 可积系统的“角变量”，但“$\la$ 的周期”给出的是作用量（action）变量，与 NRS 及本附录中作为 Darboux 坐标的长度坐标 $a_r$ 同一角色；角变量应是其共轭的扭转（twist）坐标。]
% [Tutor Note: 共形极限可直接用附录的关系式核对，这是本节最容易被跳过的一步。附录给出 $\nabla_{\zeta,R}=\zeta^{-1}R\vf+D_{\CE,h}+R\zeta\vf^{\dagger_h}$，以及 $\ep=\zeta/R$、$\nabla'_\ep=\ep\pa_{\CE,h}+\vf$。取 $R\ra0,\ \zeta\ra0$、$\zeta/R=\ep$ 固定，则 $R\zeta=\zeta^2/\ep\ra0$ 而 $\zeta^{-1}R=R/\zeta=1/\ep$，故 $\nabla_{\zeta,R}$ 的 $(1,0)$ 部分趋于 $\pa_{\CE,h}+\frac{1}{\ep}\vf$，即 $\ep\nabla'\ra\ep\pa_{\CE,h}+\vf$ 恰为 $\ep$-联络（oper 型）。同时自对偶方程 $F_{\CE,h}+R^2[\vf,\vf^{\dagger_h}]=0$ 退化为 $F_{\CE,h}=0$，即 $h$-酉联络平坦、和乐为实。这恰好解释了脚注中“$(u,0)$ 在 NAH 下的像有实和乐、只离散地与 opers 相交”：不用共形极限时 $(u,0)$ 并非 oper；只有在共形极限下 $\vf$ 的 Higgs 部分才“存活”为 oper 联络的 $u$，$(u,0)$ 的像才落到 opers 上。换言之，式\eqref{BD-KW-conf} 里的 $\rho_u$ 与 $(u,0)$ 之间必须插入“共形极限的 NAH”这一环节，不能直接用通常的 NAH。注意别把 $\ep=\zeta/R$ 与 $R/\zeta$ 弄混：$\ep$ 小时 $R/\zeta=1/\ep$ 是大的，这正是 $\frac{1}{\ep}\vf$ 这一项。]
% [Tutor Note: 关于式\eqref{BD-KW} 的三点提醒。(1) 必须限制到 $\mathrm{Bun}_G^{\rm vs}$（不含幂零 Higgs 场的“非常稳定”丛）：那里 $\De_\chi$ 的纤维是有限维的，左端 $\mathrm{Hom}_{\CM_H(G)}(\mathsf B_{\rm cc},\mathsf F_{\mu(\chi)})$ 作为两个拉格朗日 brane 的交（一般为横截的有限集）也有限维，同构才可能成立；原文自承“维数可能跳跃”，见 \cite[Section 9.5]{Fr07}——在非非常稳定丛处纤维会变成无穷维（额外的退化表示出现），故不要把式\eqref{BD-KW} 读成 $\mathrm{Bun}_G$ 上向量丛的同构。(2) “右端不依赖 $\CE$”用的是共形块丛在 $\mathrm{Bun}_G$ 上的典范平坦联络（KZ/BGG 联络），沿它平移即可等同不同 $\CE$ 处的纤维；这一步与“非常稳定”条件是两回事，勿混。(3) 第 7.6 小节中 Higgs 对的扭变可逐项核对：取 $E=K^{1/2}\oplus K^{-1/2}$，则 $\mathrm{End}(E)\ot K$ 的右上元属于 $\mathrm{Hom}(K^{-1/2},K^{1/2})\ot K=K^2$，故 $u\in H^0(C,K^2)$ 正确；左下元属于 $\mathrm{Hom}(K^{1/2},K^{-1/2})\ot K=\CO$，故常数 $1$ 正确；对角元必须为零——这正是 oper 条件（$\vf$ 关于 filtration 严格下三角）。注意 $K^{\pm1/2}$ 的存在依赖一个自旋结构（$\Theta$-特征）的选取，故“$\CL_0=\pi^*K^{1/2}$ 是典范原点”只在固定该自旋结构后成立；$SL(2)$ 与 $GL(2)$ 的差别（Prym 对 Jacobian）正是这一 $\mathbb{Z}_2$ 数据造成的，附录中“Jacobian 若 $GL(2)$、Prym 若 $SL(2)$”即指此。]
% [Tutor Note: 展望段的式\eqref{spectral} 是个直积分，其中测度 $d\mu_\si(a)$ 原文未给出；按 \cite{T10} 的谱分解，它应是复 Fenchel--Nielsen 长度坐标上的 Plancherel 型测度（含 Liouville 理论的反射振幅/DOZZ 密度），而非平凡的 Lebesgue 测度，否则“$\sigma$-无关性”无法成立。另注意：右端对每个 $\si$ 给出一个分解，而左端与 $\si$ 无关，因此“不同裤分解之间的线性关系”正是使该分解自洽的关键；其经典极限是生成函数为 \cite{NRS} 的 Yang--Yang 函数 $\mathcal Y$ 的辛变换（即 Fenchel--Nielsen 坐标变换），量子化后应是量子膨胀/量子二对数（quantum dilogarithm）型的变换。原文只说“将在别处讨论”，此处不宜过度解读为已证的定理。]
